\section{Introduction}

\subsection{Background \& Motivation}

% NOTE: Should I explain that moderns software testing approaches do not apply to DNN verifications sf
AI, particularly deep neural networks (DNNs), has gained considerable traction in recent years. Consequently, their use in high-risk domains has increased. However, it has become evident that DNNs are not free from potential errors, such as exhibiting erroneous behaviour in response to trivial perturbations in the input \cite{intriguing_properties_of_nn}. Therefore, ensuring the robustness of DNNs rigorously is essential.

DNN verifiers address the challenging task of ensuring that a neural network satisfies specific safety, robustness, or fairness properties. However, these verifiers are not free from implementation bugs. Furthermore, most state-of-the-art DNN verifiers rely on floating-point arithmetic, which results in serious scalability concerns due to numerical instability \cite{proof_minimization_in_nn_verification}. These concerns can be exploited in order to compromise a fully verified neural network \cite{exploiting_verified_nn}, thereby defeating the purpose of DNN verifiers by raising unsoundness issues.

To ensure the correctness of a DNN verifier, one could consider verifying it by formalising and then proving its implementation. Nevertheless, since DNN verifiers are very complex programs that can even rely on multiple external software components, classical program verification techniques are not feasible. For that reason, the idea of \textit{Self-Certification} \cite{program_correctness} was introduced. \textit{Self-Certification} introduces a lightweight alternative to correctness by generating and checking certificates for every program execution, rather than verifying the program as a whole.

The software responsible for generating and verifying the certifications is a \textit{proof checker}. The role that the proof checker takes in a DNN verification pipeline is shown in \autoref{fig:dnn_verification_pipeline}. The proof checker verifies the correctness of the certificates generated by the verifier when no unsafe behaviour is found, meaning that the verifier was unable to find any counterexample witnessing an erroneous behaviour of the DNN.

Despite the significant advances made in the state-of-the-art deep neural network verifier Marabou, the standard implementation of its proof checker provides no formal guarantees of its own correctness. To overcome this limitation, a fully verified implementation of Marabou's proof checker was developed in Imandra \cite{certified_proof_checker}, an industrial interactive theorem prover (ITP). In addition, Imandra supports arbitrary-precision real arithmetic, eliminating floating-point imprecision and ensuring numerical soundness.


Although this implementation demonstrates the feasibility of a certified proof checker for Marabou, its impact is limited by the adoption and maturity of Imandra within the current theorem proving landscape. More widely used ITPs, such as Rocq (formerly Coq), Lean, and Isabelle, benefit from larger user communities and more mature ecosystems \cite{survey_coq_users}. Furthermore, Imandra's comparatively limited ecosystem can increase code complexity and reduce maintainability relative to an equivalent implementation in a more established theorem prover.

Beyond ecosystem considerations, Imandra and Rocq are founded on different design philosophies. While Imandra emphasizes automation and formal reasoning through a computational logic, Rocq provides a richer type-theoretic framework for formal verification. Consequently, reimplementing the proof checker in Rocq provides an opportunity to explore the trade-offs between these approaches and assess their impact on the development of certified verification tools.

\begin{figure}[h]
\centering

\begin{tikzpicture}[
    box/.style={draw, rounded corners, align=center, minimum width=2cm, minimum height=1.0cm},
    arrow/.style={-Latex, thick}
]

\tikzset{
    verifier/.style={box, fill=blue!10},
    input/.style={box, fill=gray!10},
    sat/.style={box, fill=red!15},
    unsat/.style={box, fill=green!15},
    check/.style={box, fill=gray!15},
    finalsafe/.style={box, fill=green!25},
    finalunsafe/.style={box, fill=red!25}
}

% Center
\node[verifier] (b3) {DNN Verifier};

% Inputs (closer)
\node[input, left=1cm of b3, yshift=0.6cm] (b1) {DNN};
\node[input, left=1cm of b3, yshift=-0.6cm] (b2) {Property\\Spec.};

% Branch nodes (closer)
\node[sat, right=1cm of b3, yshift=1cm] (sat) {SAT\\(Unsafe)};
\node[unsat, right=1cm of b3, yshift=-1cm] (unsat) {UNSAT\\(Safe)};

% SAT path (compressed)
\node[check, right=1cm of sat] (sat_check) {Counterexample\\validation};

% UNSAT path (compressed)
\node[check, right=1cm of unsat] (proof) {Proof\\Checker};

% Arrows into verifier
\draw[arrow] (b1.east) -- (b3.west);
\draw[arrow] (b2.east) -- (b3.west);

% Branching
\draw[arrow] (b3.east) -- (sat.west);
\draw[arrow] (b3.east) -- (unsat.west);

% SAT flow
\draw[arrow] (sat.east) -- (sat_check.west);

% UNSAT flow
\draw[arrow] (unsat.east) -- (proof.west);

\end{tikzpicture}

\caption{The DNN verification pipeline. When the verifier finds no unsafe behaviour (UNSAT), it emits a certificate which the proof checker independently validates. When unsafe behaviour is found (SAT), a counterexample is tested directly against the DNN.}
\label{fig:dnn_verification_pipeline}
\end{figure}


\subsection{Aims and Objectives}

The primary aim of this project is to implement a certified proof checker for Marabou in Rocq, making certified DNN verification accessible within a mature, open-source ITP with a large ecosystem and community. The Imandra proof checker by Desmartin et al. \cite{certified_proof_checker} demonstrates feasibility, but its impact is limited by Imandra's comparatively small ecosystem. By reimplementing the checker in Rocq, we aim to combine DNN-specific certificate verification with the advantages of a well-established theorem prover.

The objectives of the project are as follows:

\begin{itemize}
	\item \textbf{Implement the proof checker.} Reconstruct the proof tree validation algorithm in Rocq, using the MathComp library for matrix and vector operations and Rocq's dependent types for strong invariants.

	\item \textbf{Prove soundness.} Formally establish in Rocq that any certificate accepted by the proof checker constitutes a valid proof of unsatisfiability of the corresponding DNN verification query.

	\item \textbf{Evaluate the implementation.} Compare the Rocq proof checker against the Imandra baseline in terms of functional capability and maintainability to assess the trade-offs between the two ITPs.
\end{itemize}

\subsection{Contributions}

The main contributions of this dissertation are as follows:

\begin{itemize}
	\item \textbf{A certificate parser for Marabou proof certificates.} A Haskell program that translates Marabou's JSON proof certificates into Gallina definitions. As part of the translation, the parser restructures the certificate data into a form more suitable for verification, eliminating the need for additional preprocessing within Rocq.

	\item \textbf{A proof checker for Marabou in Rocq.} The checker implements the recursive verification of Marabou's proof trees, using Rocq's dependent type system and the MathComp library. It validates case splits at internal nodes and verifies that leaf-level Farkas coefficients constitute a valid witness of unsatisfiability under the current bounds.

	\item \textbf{A formal proof of soundness in Rocq.} A Rocq theorem establishes that any certificate accepted by the proof checker genuinely certifies unsatisfiability. The proof formalises the DNN Farkas lemma adapted from the Imandra development and is fully checked by Rocq's kernel.

	\item \textbf{A comparative evaluation of the Rocq and Imandra implementations.} The two implementations are assessed across functional capability (supported certificate features) and maintainability (code size, proof complexity, library dependencies).
\end{itemize}

\subsection{Dissertation Structure}

The remainder of this dissertation is structured as follows. \autoref{sec:background} presents the necessary background on DNN verification, Marabou and interactive theorem provers. \autoref{sec:related-work} reviews related work in ITP-based neural network verification, proof production, and certified proof checking.  \autoref{sec:implementation} details the implementation of the certificate parser, proof checker, and soundness proof. \autoref{sec:evaluation} evaluates the implementation against the Imandra baseline. \autoref{sec:future-work} discusses directions for future work, and \autoref{sec:conclusion} concludes the dissertation.

